\subsection{MAS-07--09: smooth-test obstruction, heat realization and the actual metric}
Retain the fixed setup of \fid{MAS-10}: center $1/2$, strip half-width $\alpha=1/2$, trivial compact character, compact Haar mass one, and
\[
 H=L^2(\R,w(u)\,du),\qquad
 w(u)=e^{|u|}(1+u^2)^{s/2},\quad s>1,\qquad E=H/V.
\]
The full interior jets are indexed by $z_\rho=\rho-1/2$ and their orders below the zero multiplicity. The source input remains the conditional full spectral-annihilator claim of \cite[Theorem 10.2.2.1, (10.2.2.4)--(10.2.2.6), p.~446]{masullo}. The preceding Green construction has already proved $R_+\in\mathcal S_1(E_+)$ and $\det(I-wR_+)=G(w)$ from that claim, without RH or the draft's separate nuclearity theorem. The heat construction below additionally uses the latter theorem and the stipulated RH. The final no-similarity conclusion will again require only the spectral-annihilator input.

\paragraph{MAS-07: no single compact smooth test has the required multiplier.}
No compact smooth test $\phi$ has centered Mellin values $-1/z_\rho^2$ at every nontrivial zero. More generally, the same is true of $1/(c^2-z_\rho^2)$ for any fixed $c\ge0$, including the multiplier of $C$ when $c=1/2$.
\begin{proof}
Compact-character averaging and the fixed centering send $\phi$ to some $h\in C_c^\infty(\R)$, with centered Mellin transform
$F(z)=\int h(u)e^{zu}\,du$. Compact Haar mass and compact radial support justify differentiation under this average. Three integrations by parts, with vanishing boundary terms, give
\[
 |z|^3|F(z)|\le
 \sup_{|\sigma|\le1/2}\int_\R |h'''(u)|e^{\sigma u}\,du<\infty
 \qquad(z\ne0,\ |\Re z|\le1/2).
\]
Classically $|\Re z_\rho|<1/2$, and there are zeros with $|z_\rho|\to\infty$. But
$|z_\rho|^3/|c^2-z_\rho^2|\sim|z_\rho|\to\infty$, a contradiction. The denominators do not vanish at a nontrivial zero, since such a zero is nonreal. This proof needs no RH. Under RH it is exactly the obstruction to the prescribed values $1/(\gamma^2+c^2)$ at $z_\rho=i\gamma$.
\end{proof}
The assertion concerns the designated zero-by-zero multiplier, not an unordered eigenvalue multiset after an unspecified permutation. A different kernel or subsequent functional calculus is not excluded.

\paragraph{MAS-08: trace-class heat operators and the same inverse square.}
For this paragraph invoke the stipulated RH and write $\rho=1/2+i\gamma$, $\gamma\in\R\setminus\{0\}$, with multiplicity $N_\gamma$. Every zero is strictly interior to the fixed strip, so all jets
$\ell_{\gamma,m}(f)=\int f(u)e^{i\gamma u}u^m\,du$, $0\le m<N_\gamma$, are retained. In addition to their finite independence and norm-dense span in $E'$, assume the draft's continuous map
\[
 \mathcal A:C_c^\infty(\R)\longrightarrow\mathcal S_1(E)
\]
obtained from the centered radial sector of \cite[Theorem 10.3.2.1(2), (10.3.2.2), pp.~453--454]{masullo}. Its continuity is in trace norm, not just operator norm.

Let $T_vf(u)=f(u-v)$. The proof in \fid{MAS-10} gives a strongly continuous quotient group with
$\|T_v\|\le e^{B_0|v|}$, $B_0=1/2+s/4$. For any $h$ satisfying
$\int|h(v)|e^{B_0|v|}\,dv<\infty$, define
\[
 \mathcal A(h)f=\int_\R h(v)T_vf\,dv,\qquad
 \|\mathcal A(h)\|\le\int_\R|h(v)|e^{B_0|v|}\,dv.
\]
This is a Hilbert-space Bochner integral for each fixed vector $f$, hence a strong operator integral. No Bochner measurability of $v\mapsto T_v$ in $\mathcal B(E)$ is asserted.
Put
\[
 k_t(v)=(4\pi t)^{-1/2}e^{-v^2/(4t)},\qquad H_t=\mathcal A(k_t)\quad(t>0).
\]
Then $H_t\in\mathcal S_1(E)$, with trace-norm continuity on compact subintervals of $t>0$, and
\[
 H_tH_r=H_{t+r},\qquad H_t\longrightarrow I\ \hbox{strongly as }t\downarrow0,
 \qquad \sup_{0<t\le1}\|H_t\|<\infty.
\]
\begin{proof}
Choose $\chi\in C_c^\infty((-1,1))$ with
$\sum_{n\in\Z}\chi(v-n)=1$. One construction divides a nonnegative bump positive on $[-1/2,1/2]$ by the positive periodic sum of its translates. Put
$q_{n,t}(v)=\chi(v)k_t(v+n)$. Then
$k_t(v)=\sum_n q_{n,t}(v-n)$, pointwise and in $L^1(e^{B_0|v|}\,dv)$.

Continuity of the linear map $\mathcal A$ on the fixed-support Fr\'echet space with support in $[-1,1]$ supplies $C,M$ such that
\[
 \|\mathcal A(q)\|_1\le C\sum_{j=0}^{M}\|q^{(j)}\|_\infty.
\]
For $t$ in any fixed interval $[a,b]\subset(0,\infty)$, Gaussian differentiation gives
\[
 \sum_{j=0}^{M}\|q_{n,t}^{(j)}\|_\infty
 \le C_{a,b}(1+|n|)^M
       \exp\!\left(-\frac{((|n|-1)_+)^2}{4b}\right).
\]
Each derivative is a polynomial times the Gaussian; on $|v|\le1$ this is the displayed bound. Since
$\mathcal A(q(\,\cdot-n))=T_n\mathcal A(q)$, multiplication by $e^{B_0|n|}$ still leaves a summable bound. Thus
$\sum_nT_n\mathcal A(q_{n,t})$ converges in trace norm, uniformly on $[a,b]$. The weighted $L^1$ estimate identifies its sum with the strong integral $H_t$. Each summand is trace-norm continuous in $t$, by the conditional test-map continuity, so their sum is too.

For each fixed $f$, vector-valued Fubini and the group law give
\[
 H_tH_rf=\iint k_t(v)k_r(y)T_{v+y}f\,dy\,dv=H_{t+r}f.
\]
The integrable majorant is
$k_t(v)k_r(y)e^{B_0(|v|+|y|)}\|f\|$; Gaussian convolution gives the last equality. This is vectorwise Fubini, not an operator-valued Bochner assertion about the translation group.
Finally substitute $v=\sqrt t\,y$ in the norm bound. For $0<t\le1$ the integrand is dominated by a fixed Gaussian times $e^{B_0|y|}$. This proves uniform boundedness. Applying the same substitution to
$H_tf-f$ and using strong continuity of $T_v$ proves strong convergence by dominated convergence.
\end{proof}

The complete spectral action, including jets, is
\begin{equation}\label{eq:MasHeatJets}
 H_t'\ell_{\gamma,m}
 =\sum_{j=0}^m\binom mj
 F_t^{(m-j)}(i\gamma)\ell_{\gamma,j},
 \qquad F_t(z)=e^{tz^2}.
\end{equation}
Consequently the nonzero spectrum of $H_t$ consists precisely of $e^{-t\gamma^2}$, with algebraic multiplicity $2N_\gamma$ for each $\gamma>0$.
\begin{proof}
Apply a jet to the strong integral and expand $(u+v)^m$ in the translation formula. Gaussian decay dominates every polynomial moment and the exponential factor. The resulting moments are derivatives of the Gaussian Mellin transform $F_t$, proving \eqref{eq:MasHeatJets}. These finite jet actions obey the product rule $M_FM_G=M_{FG}$ by Leibniz's identity.

If a nonzero eigenvalue $\lambda$ of the compact $H_t$ differs from all displayed values, evaluation on an eigenvector by jets of order zero and then successively higher orders makes every evaluation zero. Their totality makes that vector zero, a contradiction. Conversely an order-zero eigenfunctional at a displayed value makes $H_t-\lambda I$ noninvertible, so that value is in the spectrum.

For a displayed isolated value $\lambda$, take its finite-rank Riesz projection. Its transpose acts on a finite jet block as identity when the block's diagonal value is $\lambda$ and as zero otherwise: restrict the resolvent to that invariant block and integrate its finite triangular matrix. Exactly the two blocks at $\gamma$ and $-\gamma$ match, with total dimension $2N_\gamma$. The image of the dense jet span is their finite-dimensional span, which is closed; continuity makes the entire projection range equal that span. This gives the full algebraic multiplicity. These are generalized eigenspaces, or Riesz subspaces; no diagonalizability is presumed at a multiple zero.
\end{proof}

The even kernels commute with the unitary involution $Jf(u)=f(-u)$ from \fid{MAS-10}. Hence $H_{t,+}=H_t|_{E_+}$ is trace class. On each paired Riesz subspace the dual even part has the independent basis
\[
 \ell_{\gamma,m}+(-1)^m\ell_{-\gamma,m},
 \qquad 0\le m<N_\gamma.
\]
Indeed $J'\ell_{\gamma,m}=(-1)^m\ell_{-\gamma,m}$ interchanges the two blocks, and applying $(I+J')/2$ to the dense jet span proves totality on the even part. Thus the algebraic multiplicity of $e^{-t\gamma^2}$ on $E_+$ is $N_\gamma$.

The heat construction gives the exact identity with the Green operator:
\begin{equation}\label{eq:MasHeatInverse}
 Q=\int_0^1H_t\,dt,\qquad
 R=Q(I-H_1)^{-1},\qquad
 R_+=Q|_{E_+}(I-H_{1,+})^{-1}.
\end{equation}
\begin{proof}
Here $Q$ really is a $\mathcal B(E)$-valued Bochner integral: $t\mapsto H_t$ is norm continuous for $t>0$, has separable norm range there, and is bounded near zero. Assigning any value at $0$ does not change the integral. Its truncations $\int_\varepsilon^1H_tdt$ are trace class, and the omitted interval has operator norm at most a constant times $\varepsilon$, so $Q$ is compact. Since $\gamma\ne0$, the just-proved spectrum excludes $1$ from $\operatorname{Spec}H_1$; hence $I-H_1$ has a bounded inverse.

On every finite jet block the operator $Q(I-H_1)^{-1}$ has analytic multiplier
\[
 \frac{\int_0^1e^{tz^2}\,dt}{1-e^{z^2}}=-\frac1{z^2}
 \quad\hbox{near }z=i\gamma.
\]
The denominator is nonzero there. Integration and this finite jet algebra identity include every derivative up to the full zero multiplicity. This is exactly the multiplier of the bounded Green operator $R$ in \fid{MAS-10}. Their transposes therefore agree on all jets, whose span is dense in $E'$, so the bounded operators agree. Restricting their commuting even parts proves the last identity.
\end{proof}
The trace-class property and ordinary determinant of this $R_+$ were already established by \fid{MAS-10}; the identity also gives $Q=R(I-H_1)\in\mathcal S_1(E)$. No estimate of $\int_0^1\|H_{t,+}\|_1dt$ is needed for these conclusions. For comparison, the scalar trace integral is finite:
\[
 \int_0^1\Tr H_{t,+}\,dt
 =\sum_{\gamma>0}N_\gamma\frac{1-e^{-\gamma^2}}{\gamma^2}<\infty.
\]
At each positive time the Hilbert trace theorem and the proved multiplicities give the trace; Tonelli and the zero count give the integral. Scalar trace integrability alone would not justify a trace-norm integral for these nonpositive operators.

\paragraph{MAS-09: the actual quotient metric rules out selfadjointness.}
Under stipulated RH define
\[
 K(x)=\int_\R\frac{e^{ixu}}{w(u)}\,du.
\]
This is a continuous real even function, tends to zero as $|x|\to\infty$, and satisfies $K(x)>0$ for every real $x$.
\begin{proof}
The reciprocal weight is integrable, giving continuity and decay by the $L^1$ Fourier theorem. Put $\beta=s/2>0$. The Gamma integral gives
\[
 (1+u^2)^{-\beta}
 =\frac1{\Gamma(\beta)}\int_0^\infty r^{\beta-1}e^{-r}e^{-ru^2}\,dr.
\]
For $r>0$, use the Cauchy Fourier identity
$e^{-|u|}=\pi^{-1}\int_\R e^{iyu}(1+y^2)^{-1}dy$. For $u\ne0$ it follows by a single residue at $i$ or $-i$ in the decaying half-plane; at $u=0$ it is the arctangent integral. The Gaussian Fourier integral now gives
\[
 \int_\R e^{ixu}e^{-|u|}e^{-ru^2}\,du
 =\frac1\pi\sqrt{\frac\pi r}
   \int_\R\frac{e^{-(x+y)^2/(4r)}}{1+y^2}\,dy>0.
\]
Fubini here is justified by the integrable Cauchy density and Gaussian. Integrate this strictly positive expression against the positive Gamma measure. Changing the original $r,u$ order is justified by the absolute integrable majorant
$e^{-|u|}(1+u^2)^{-\beta}$. This proves strict positivity. Reality and evenness also follow directly from the even reciprocal weight.
\end{proof}

Use the paper's conjugate-linear-first inner-product convention, so that the actual Riesz representation is $\ell_{\gamma,0}(f)=\langle v_\gamma,f\rangle$, where
$v_\gamma(u)=e^{-i\gamma u}/w(u)$. Since the functional annihilates $V$, this vector belongs to $V^\perp$ and represents it isometrically on the Hilbert quotient. In $E_+$ the even representative is
$c_\gamma(u)=\cos(\gamma u)/w(u)$, and direct integration gives
\begin{equation}\label{eq:MasActualCorrelation}
 \langle c_\gamma,c_\delta\rangle
 =\tfrac12\bigl(K(\gamma-\delta)+K(\gamma+\delta)\bigr)>0.
\end{equation}
But
\[
 R_+^*c_\gamma=\gamma^{-2}c_\gamma,\qquad
 H_{t,+}^*c_\gamma=e^{-t\gamma^2}c_\gamma.
\]
Two distinct positive ordinates give distinct real eigenvalues and nonorthogonal eigenvectors. A selfadjoint operator has orthogonal eigenvectors at distinct real eigenvalues, so neither $R_+$ nor any $H_{t,+}$ is selfadjoint or positive in this metric. More generally, this same argument applies to any bounded multiplier whose real values differ at two positive ordinates. This is a calculation in the specified quotient norm, not a replacement model.

\paragraph{MAS-09: no bounded similarity can repair $R_+$.}
Conditional only on the full spectral-annihilator input of \fid{MAS-10}, $R_+$ is not boundedly similar to a selfadjoint operator.
\begin{proof}
First, if RH fails, choose an off-line nontrivial zero. Its centered parameter $z=\rho-1/2$ has nonzero real part and nonzero imaginary part, since zeta has no real nontrivial zero. Therefore $z^2$, and hence $-1/z^2$, is nonreal. By \fid{MAS-10} this is an eigenvalue of $R_+$. Similarity preserves spectrum, whereas a selfadjoint operator has real spectrum, which is already a contradiction.

It remains to consider RH. If some zero has $N_\gamma\ge2$, the independent even jets
$e_0=\ell_{\gamma,0}+\ell_{-\gamma,0}$ and
$e_1=\ell_{\gamma,1}-\ell_{-\gamma,1}$ satisfy, with $F(z)=-1/z^2$,
\[
 R_+'e_0=\gamma^{-2}e_0,\qquad
 R_+'e_1=\gamma^{-2}e_1+F'(i\gamma)e_0,\qquad
 F'(i\gamma)=\frac{2}{(i\gamma)^3}\ne0.
\]
The full interior jets are essential here. Through the Riesz identification this gives a nontrivial generalized eigenvector for $R_+^*$ at a real eigenvalue. A selfadjoint operator $A$ has none: if
$(A-\lambda)y=x$ and $(A-\lambda)x=0$ with $\lambda\in\R$, then
$\|x\|^2=\langle x,(A-\lambda)y\rangle=\langle(A-\lambda)x,y\rangle=0$.
Bounded similarity preserves this obstruction.

If all zeros are simple, the positive zero count $N(T)\sim T\log T/(2\pi)$ forces distinct positive pairs $\gamma_n,\delta_n\to\infty$ with $|\gamma_n-\delta_n|\to0$. Otherwise a positive eventual lower gap would give $N(T)=O(T)$; equivalently, apply the pigeonhole principle to the successive gaps below $T$. Discreteness ensures pairs with gaps tending to zero escape each fixed compact interval.
By \eqref{eq:MasActualCorrelation},
\[
 \|c_\gamma\|^2=\tfrac12(K(0)+K(2\gamma))\longrightarrow K(0)/2,
 \qquad
 \langle c_{\gamma_n},c_{\delta_n}\rangle\longrightarrow K(0)/2.
\]
Thus their normalized vectors $u_n,v_n$ satisfy
$\langle u_n,v_n\rangle\to1$ and $\|u_n-v_n\|\to0$, while remaining eigenvectors of $R_+^*$ for distinct eigenvalues.

If $R_+^*=SAS^{-1}$ with $A$ selfadjoint and $S$ boundedly invertible, then
$S^{-1}u_n,S^{-1}v_n$ are orthogonal, each of norm at least $1/\|S\|$. Their difference has norm at least $\sqrt2/\|S\|$, contradicting its upper bound
$\|S^{-1}\|\|u_n-v_n\|\to0$. Finally, bounded similarity of $R_+$ to a selfadjoint operator would give one for $R_+^*$ by taking adjoints. This proves the assertion in all cases, without a simplicity assumption.
\end{proof}

Consequently the actual ordinary Fredholm determinant constructed in \fid{MAS-10} can equal that of the designated positive diagonal $B$ while its quotient operator has no bounded similarity to a selfadjoint operator. The obstruction distinguishes the determinant and full algebraic spectrum from the Hilbert metric. A change to an unbounded or nonequivalent completion would be another construction; it is not supplied by this argument. Every Masullo-dependent conclusion remains conditional on the exact draft inputs stated above, with nuclearity needed for the Gaussian trace-class family but not for \fid{MAS-10} or the no-similarity conclusion.
